\chapter{I dati: dal mondo reale al dataset}\label{cap:dati}

% Capitolo 3 — I dati: dal mondo reale al dataset.

Ogni progetto di machine learning comincia dai dati. In questo capitolo
vediamo le sei feature cliniche del progetto, perché il dataset è sintetico,
come è stato diviso e congelato, come viene normalizzato e perché tutto
questo rende il progetto riproducibile \cite{progetto2026}. Le etichette del
dataset le mette il teacher rule-based, che è l'argomento del capitolo
\ref{cap:teacher}.

\section{Le sei feature cliniche}\label{sec:sei-feature-cliniche}

Le sei feature sono parametri vitali misurati su un paziente. L'ordine
ufficiale è fissato dal contratto in \codice{FEATURE\_ORDER} dentro
\file{src/telemedicina\_supervised/safety/safety\_rules.py}: è l'ordine del
vettore in ingresso al modello e non va modificato senza aggiornare il
contratto.

La tabella seguente riporta, per ogni feature, l'unità di misura, il range
normale e un breve significato clinico. I range normali sono quelli reali
del progetto, definiti in \codice{RANGE\_NORMALI} in \file{safety\_rules.py}.

\begin{table}[h]
\centering
\caption{Le sei feature cliniche: unità, range normale e significato.}
\label{tab:feature-cliniche}
\begin{tabular}{llll}
\toprule
\textbf{Feature} & \textbf{Unità} & \textbf{Range normale} & \textbf{Significato clinico} \\
\midrule
\codice{pressione\_sistolica} & mmHg & 90--140 & pressione massima del cuore \\
\codice{pressione\_diastolica} & mmHg & 60--90 & pressione minima tra i battiti \\
\codice{frequenza\_cardiaca} & bpm & 60--100 & battiti al minuto \\
\codice{temperatura} & °C & 36.0--37.5 & temperatura corporea \\
\codice{saturazione\_ossigeno} & \% & 95--100 & ossigeno nel sangue \\
\codice{glicemia} & mg/dL & 70--140 & zucchero nel sangue \\
\bottomrule
\end{tabular}
\end{table}

Oltre ai range normali, \file{safety\_rules.py} definisce i \textbf{range
fisiologici} (i valori "possibili" per un essere umano) e le \textbf{soglie
critiche e moderate} che determinano la classe di rischio. Un input con un
valore fuori dal range fisiologico non è un paziente a rischio basso: è una
misurazione invalida, e il sistema risponde \classe{errore}, mai una classe
di rischio. Questo è un dettaglio importante: il modello lavora solo su
input validi.

\section{Perché un dataset sintetico}\label{sec:perche-dataset-sintetico}

Il dataset del progetto è \textbf{sintetico}: non contiene dati reali di
pazienti, ma campioni generati artificialmente. Le ragioni sono tre.

\begin{enumerate}
  \item \textbf{Privacy}: i dati sanitari reali sono estremamente sensibili.
        Un dataset sintetico elimina ogni problema di consenso e
        anonimizzazione.
  \item \textbf{Controllo della distribuzione}: possiamo decidere quanti
        campioni generare per ogni profilo clinico, assicurandoci che tutte
        le classi siano rappresentate e che le regioni anomale siano
        sovracampionate.
  \item \textbf{Riproducibilità}: con un seed fisso, la generazione produce
        sempre gli stessi dati. Stessi dati, stessi numeri, sempre.
\end{enumerate}

Come funziona la generazione? Il generatore
(\file{training/synthetic\_generator.py}) usa una \emph{mixture di profili}:
ogni profilo rappresenta una condizione clinica (normale, ipertensione,
tachicardia, febbre, ipossia, ipoglicemia, shock, ecc.) con un peso e, per
ogni feature, una coppia (media, deviazione standard) di una distribuzione
gaussiana. Il campionamento avviene in tre passi:

\begin{enumerate}
  \item si sceglie un profilo con probabilità proporzionale al suo peso;
  \item si campionano i sei valori dalle gaussiane del profilo, limitandoli
        ai range fisiologici;
  \item si impone il vincolo strutturale sistolica $>$ diastolica (con
        differenza minima di 10 mmHg), ricampionando la diastolica se
        necessario.
\end{enumerate}

C'è poi un dettaglio importante: la \textbf{buffer zone}. Vicino alle soglie
di classificazione, un piccolo rumore di misura può far cambiare classe a un
campione, creando etichette "rumorose" sul gradino. Per questo il generatore
scarta i campioni che cadono entro un margine da una qualunque soglia di
confine (critica, moderata o limite del range normale). I margini reali sono
in \codice{MARGINI\_BUFFER}: 2.0 mmHg per le pressioni, 1.0 bpm per la
frequenza, 0.2 °C per la temperatura, 1.0\% per la saturazione e 2.0 mg/dL
per la glicemia.

\begin{esempio}{Quanti campioni vengono scartati}
La buffer zone scarta parecchio: per generare i 40\,000 campioni di train
sono stati scartati 52\,504 campioni; per la validation 13\,192 su 10\,000
richiesti e per il test 26\,341 su 20\,000 (da
\file{data/processed/metadati.json}). È il prezzo da pagare per avere
etichette pulite vicino alle soglie.
\end{esempio}

Ogni campione generato viene etichettato dal teacher rule-based
(capitolo \ref{cap:teacher}): la label è la classe di rischio calcolata
dalle regole. I campioni che il teacher giudica \classe{errore} (input
invalidi) vengono scartati: il dataset di training non contiene mai la
classe \classe{errore}.

\section{Lo split congelato}\label{sec:split-congelato}

Il dataset è diviso in tre blocchi, generati con seed \textbf{derivati e
diversi}:

\begin{table}[h]
\centering
\caption{Lo split congelato del dataset (seed base 41).}
\label{tab:split}
\begin{tabular}{lrr}
\toprule
\textbf{Blocco} & \textbf{Seed} & \textbf{Campioni} \\
\midrule
train & 41 & 40\,000 \\
validation & 42 & 10\,000 \\
test & 43 & 20\,000 \\
\midrule
totale & & 70\,000 \\
\bottomrule
\end{tabular}
\end{table}

Lo split è \textbf{congelato}: i tre blocchi vengono generati una volta e
salvati su disco in \file{data/processed/} (\codice{X\_train.npy},
\codice{y\_train.npy}, ecc.). Durante il tuning degli iperparametri si
consultano solo train e validation; il test resta intatto fino alla
valutazione finale. Se il test venisse usato per scegliere la configurazione
migliore, le metriche finali sarebbero ottimistiche: il modello avrebbe
"visto" il giudice prima dell'esame.

\begin{esempio}{Il test non si tocca mai}
In \file{training/train\_model.py} la selezione della configurazione migliore
avviene sulla \emph{validation loss} (\codice{seleziona\_migliore}); il test
viene usato una sola volta, alla fine, per calcolare le metriche finali.
\end{esempio}

\section{La normalizzazione}\label{sec:normalizzazione}

Le sei feature hanno unità di misura molto diverse: la pressione sistolica
si misura in decine di mmHg, la temperatura in gradi attorno a 37, la
saturazione in percentuale tra 95 e 100. Se le passassimo al modello così
come sono, le feature con valori più grandi dominerebbero il calcolo delle
distanze e dei gradienti, indipendentemente dalla loro importanza clinica.

La soluzione è la \textbf{standardizzazione}: ogni feature viene trasformata
in

\begin{equation}
x' = \frac{x - \mu}{\sigma}
\end{equation}

dove $\mu$ è la media e $\sigma$ la deviazione standard della feature,
calcolate \textbf{solo sul train}. Dopo la trasformazione ogni feature ha
media 0 e deviazione standard 1, e tutte le scale sono confrontabili.

Nel progetto la standardizzazione è implementata in
\file{src/telemedicina\_supervised/ml/scaler.py} (\codice{StandardScaler}):
il metodo \codice{fit} stima media e deviazione sul train, \codice{transform}
applica la trasformazione. Lo scaler viene salvato \emph{dentro} l'artifact
del modello (\file{data/models/mlp\_telemedicina.npz}), così a runtime
l'agente applica esattamente la stessa trasformazione vista in training.

\begin{esempio}{Perché lo scaler si calcola solo sul train}
Se media e deviazione fossero calcolate su tutto il dataset (test incluso),
il test "perderebbe" la sua natura di dati mai visti: è una forma di
\emph{leakage}. Il codice di \file{training/train\_model.py} è esplicito:
\codice{scaler = StandardScaler().fit(dati["X\_train"])}.
\end{esempio}

\section{Perché i dati congelati contano}\label{sec:dati-congelati-riproducibilita}

Congelare i dati non è un dettaglio burocratico: è ciò che rende il progetto
\textbf{riproducibile}. Con lo stesso seed base (41) e lo stesso codice, la
generazione produce file byte-identici: chiunque riesegua la pipeline ottiene
gli stessi numeri, dallo split alle metriche finali.

Per rendere la provenienza verificabile, \file{data/processed/metadati.json}
registra gli \textbf{hash di provenienza} (sha256, primi 12 caratteri) dei
file che determinano i dati:

\begin{itemize}
  \item \codice{hash\_safety\_rules}: l'hash di \file{safety\_rules.py}, la
        fonte unica delle soglie;
  \item \codice{hash\_synthetic\_generator}: l'hash del generatore;
  \item \codice{hash\_teacher\_rules}: l'hash del teacher.
\end{itemize}

Se una soglia clinica cambia, l'hash cambia e i dati generati con la vecchia
versione non sono più confrontabili con i nuovi. È un modo semplice ma
efficace per sapere \emph{quale versione del codice ha prodotto quali dati}.
Il tema della riproducibilità è approfondito nel capitolo
\ref{cap:riproducibilita}.

Nel prossimo capitolo (\ref{cap:teacher}) vediamo chi produce le etichette:
il teacher rule-based.